\documentclass[../DoAn.tex]{subfiles}
\begin{document}

\section{Đặt vấn đề}
\label{sec:dvd}
\indent Nhận dạng chữ số viết tay là bước đầu tiên của các hệ thống OCR (Optical character recognition – nhận diện ký tự quang học) trong ngân hàng, bưu điện. Tuy nhiên, sự khác biệt trong nét vẽ, độ nghiêng của nét chữ số, hay độ nhiễu trong hình ảnh khiến việc nhận diện không đơn giản. 

\section{Mục tiêu và định hướng giải pháp}
Để giải quyết vấn đề này, ta đến với bài toán phân loại đa lớp. Trong học máy, đây là bài toán phân loại các mẫu dữ liệu vào một trong nhiều hơn hai lớp đối tượng, với nhiều hướng giải quyết, bao gồm các phương pháp thiên về học máy truyền thống, và các phương pháp thiên về học sâu. Với mục đích hiểu sâu hơn về cả học máy truyền thống và học sâu, em đã lựa chọn giải quyết bài toán bằng hai phương pháp: Support vector machine (SVM) và Convolutional neural network (CNN), lần lượt là các phương pháp được ứng dụng phổ biến trong học máy truyền thống và học sâu. Về bộ dữ liệu dành cho huấn luyện và kiểm tra, dataset MNIST, với 70000 ảnh chữ số viết tay đen trắng kích thước 24x24 pixel, là một lựa chọn phù hợp cho bài toán. 

\section{Bố cục đồ án}

Báo cáo được chia thành 6 chương như sau:

Chương 2 trình bày cơ sở lý thuyết của hai phương pháp chính được sử dụng: Support Vector Machine (SVM) và Convolutional Neural Network (CNN), bao gồm các khái niệm, công thức toán học và nguyên lý hoạt động.

Chương 3 mô tả quá trình áp dụng hai mô hình vào bài toán phân loại chữ số MNIST, từ tiền xử lý dữ liệu, thiết kế kiến trúc mạng, đến các tham số huấn luyện và triển khai code.

Chương 4 trình bày kết quả thực nghiệm, bao gồm độ chính xác, thời gian huấn luyện, ma trận nhầm lẫn và phân tích sai số của cả hai mô hình, đồng thời so sánh hiệu năng giữa SVM và CNN.

Chương 5 minh họa giao diện demo để kiểm tra mô hình trên các mẫu dữ liệu thực tế, cho thấy khả năng ứng dụng của hệ thống.

Chương 6 tổng kết các kết quả đạt được.